\chapter{联络、挠率与曲率}

Lie 导数借助一个给定的流比较张量，但它不能回答另一类问题：如果没有指定流，怎样把一个切向量沿任意曲线搬到邻点，并定义“平行”与“加速度”？这需要额外的几何结构——联络。联络不是由光滑流形自动给出的；它规定不同切空间之间如何做无穷小比较。挠率与曲率随后不是独立添加的公式，而是联络的两个张量性缺陷：挠率比较协变微分的反对称部分与 Lie 括号，曲率比较沿两个方向连续平行移动的不交换性。

\section{向量丛联络与 affine connection}
先给出比切丛更一般的母结构。设 $\pi:E\to M$ 为实或复向量丛。一个线性联络可以定义为
\begin{equation*}\nabla:\Gamma(E)\longrightarrow\Omega^1(M;E),
\qquad
\nabla(fs)=\dd f\otimes s+f\nabla s.
\end{equation*}
对向量场 $X$ 取收缩便得到 $\nabla_Xs:=(\nabla s)(X)$。在局部标架 $\{e_a\}$ 下，$\nabla e_b=\omega^a{}_b\otimes e_a$，因此任意截面 $s=s^ae_a$ 满足
\[
\nabla s=(\dd s^a+\omega^a{}_b s^b)\otimes e_a.
\]
联络的定义只依赖向量丛结构；切丛联络、法丛联络以及后面出现的主丛伴随丛联络都属于同一框架。

向量丛联络还可以从 principal bundle 的水平分布统一产生。设 $P\to M$ 是 principal $G$-bundle，$\mathfrak g$ 是 $G$ 的 Lie algebra。对 $\xi\in\mathfrak g$，右作用生成 fundamental vertical vector field
\[
\xi_P(p)
:=\left.\frac{\dd}{\dd t}\right|_0p\exp(t\xi).
\]
垂直子空间为
\[
V_pP=\ker\pi_{*p}
=\{\xi_P(p):\xi\in\mathfrak g\}.
\]
principal connection 可以等价地指定一个 $G$-equivariant 水平补空间
\begin{equation*}T_pP=H_pP\oplus V_pP,
\qquad
(R_g)_*H_pP=H_{pg}P,\end{equation*}
或者指定一个 $\mathfrak g$-值一形式 $\boldsymbol\omega\in\Omega^1(P;\mathfrak g)$，满足
\begin{equation}
\boldsymbol\omega(\xi_P)=\xi,
\qquad
R_g^*\boldsymbol\omega
=\operatorname{Ad}_{g^{-1}}\boldsymbol\omega.
\label{eq:principal-connection-form}
\end{equation}
此时
\[
H_pP=\ker\boldsymbol\omega_p.
\]
因此联络的几何本体是“怎样在总空间中选择水平提升”，局部矩阵值一形式只是选择截面后的表示。

若 $s_\alpha:U_\alpha\to P$ 是局部截面，定义局部 gauge potential
\[
A_\alpha=s_\alpha^*\boldsymbol\omega.
\]
在 $s_\beta=s_\alpha g_{\alpha\beta}$ 下，利用\eqnrefs{eq:principal-connection-form} 可得
\begin{equation}
A_\beta
=g_{\alpha\beta}^{-1}A_\alpha g_{\alpha\beta}
+g_{\alpha\beta}^{-1}\dd g_{\alpha\beta}.
\label{eq:principal-local-connection-transform}
\end{equation}
这就是\eqnrefs{eq:math-dg-connection-gauge-transform} 的 principal-bundle 来源。

给定表示 $\rho:G\to GL(W)$ 与 associated bundle $E=P\times_\rho W$，principal connection 唯一诱导 $E$ 上的向量丛联络。局部上若截面写成 $W$-值函数 $\psi_\alpha$，则
\begin{equation*}D\psi_\alpha
=\dd\psi_\alpha+\rho_*(A_\alpha)\psi_\alpha.\end{equation*}
换 gauge 后，$D\psi$ 按与 $\psi$ 相同的表示齐次变换。于是“非齐次变换的 connection + 齐次变换的 matter field”并不是物理中特有的记号，而是 associated bundle 的一般微分几何机制。

\begin{derivation}[局部 connection form 的非齐次 gauge 变换]
换局部截面 $s'=sg$ 时，新截面沿 $X$ 的切向量 $s'_*X$ 既含旧截面切向量的右平移，又含 $g$ 自身变化产生的垂直分量。Connection form 对水平部分用 $G$-equivariance、对垂直部分用 reproduction property，两项合起来给出非齐次变换律。

分解新截面切向量。令 $s'=sg$，其中 $g:U\to G$。对 $X\in T_xM$，先计算新截面沿 $X$ 的切向量 $s'_*X$。由乘积截面微分，$s'_*X$ 同时含有 $s_*X$ 经右平移到 $s(x)g(x)$ 的水平部分，以及 $g(x)$ 自身沿 $X$ 变化产生的垂直部分。把后者经右平移回到单位元再经 fundamental vertical field 嵌入，得
\begin{equation*}
s'_*X
=(R_g)_*s_*X+\bigl(g^{-1}\dd g(X)\bigr)_P(sg),
\end{equation*}
其中 $(R_g)_*$ 是右平移 $R_g:h\mapsto hg$ 的切映射，$\bigl(g^{-1}\dd g(X)\bigr)_P$ 是把 $\mathfrak g$ 元素 $g^{-1}\dd g(X)$ 经 fundamental field 映射 $\mathfrak g\to T_{sg}P$ 嵌入所得垂直向量。

代入 connection form。由 $A'(X)=\boldsymbol\omega_{s'(x)}(s'_*X)$ 且 $s'(x)=s(x)g(x)$，代入上式并利用 connection form 的线性性，
\begin{equation*}
A'(X)
=\boldsymbol\omega_{sg}\bigl((R_g)_*s_*X\bigr)
 +\boldsymbol\omega_{sg}\bigl(\bigl(g^{-1}\dd g(X)\bigr)_P(sg)\bigr).
\end{equation*}

处理水平项。对第一项用 $G$-equivariance $\boldsymbol\omega_{sg}\bigl((R_g)_*v\bigr)=\operatorname{Ad}_{g^{-1}}\boldsymbol\omega_s(v)$，取 $v=s_*X$，
\begin{equation*}
\boldsymbol\omega_{sg}\bigl((R_g)_*s_*X\bigr)
=\operatorname{Ad}_{g^{-1}}\boldsymbol\omega_s(s_*X)
=\operatorname{Ad}_{g^{-1}}A(X).
\end{equation*}
Equivariance 是 connection form 的定义公理之一，几何意义是"右作用把 connection form 按 adjoint 表示拉回"。

处理垂直项。对第二项用 reproduction property $\boldsymbol\omega_p\bigl(\xi_P(p)\bigr)=\xi$，取 $\xi=g^{-1}\dd g(X)$、$p=sg$，
\begin{equation*}
\boldsymbol\omega_{sg}\bigl(\bigl(g^{-1}\dd g(X)\bigr)_P(sg)\bigr)
=g^{-1}\dd g(X).
\end{equation*}
Reproduction property 是 connection form 的另一条定义公理：在 fundamental vertical field 上 connection form 是恒等映射。

合并两项。
\begin{equation*}
A'(X)
=\operatorname{Ad}_{g^{-1}}A(X)+g^{-1}\dd g(X).
\end{equation*}
对矩阵 Lie 群 $\operatorname{Ad}_{g^{-1}}A=g^{-1}Ag$，于是
\begin{equation*}
A'(X)=g^{-1}A(X)g+g^{-1}\dd g(X),
\end{equation*}
即\eqnrefs{eq:principal-local-connection-transform}。非齐次项来自"局部截面本身也在变化"，所以任何把 connection coefficient 当作普通 tensor 的做法都会漏掉这一项。

典型例子。
\begin{itemize}
\item $G=U(1)$：$\operatorname{Ad}_{g^{-1}}=\mathrm{id}$（Abel 群 adjoint 平凡），$A'=A+\ii\,\dd\chi$（取 $g=e^{\ii\chi}$）。这正是电磁势的规范变换 $A_\mu\mapsto A_\mu+\partial_\mu\chi$。
\item $G=SU(N)$：$A'=g^{-1}Ag+g^{-1}\dd g$，非 Abel 规范变换。$g^{-1}Ag$ 是共轭转动，$g^{-1}\dd g$ 是纯规范补偿。在 Yang--Mills 理论中，这一变换律保证 $D\psi=\dd\psi+A\psi$ 与 $\psi$ 同步齐次变换 $\psi\mapsto g^{-1}\psi$。
\end{itemize}

非齐次项 $g^{-1}\dd g$ 是规范势在规范变换下的"补偿"，对应物理上规范势 $A_\mu$ 在波函数相位重定向下获得的导数项。它使规范势不能视为张量，但保证协变导数 $D\psi$ 是张量（齐次变换）。这一非齐次性是规范理论全部相互作用结构的根源。
\end{derivation}
取 $E=TM$ 时，流形 $M$ 上的 affine connection 等价地写成双线性映射
\[
\nabla:\Gamma(TM)\times\Gamma(TM)\to\Gamma(TM),
\qquad
(X,Y)\mapsto\nabla_XY,
\]
满足对 $f\in C^\infty(M)$
\begin{align*}
\nabla_{fX}Y&=f\nabla_XY,\\
\nabla_X(fY)&=X(f)Y+f\nabla_XY.
\end{align*}
第一条说明它对“求导方向”是张量性的；第二条说明它对被求导向量满足 Leibniz 规则。与 Lie 导数不同，$\nabla_XY$ 对交换 $X,Y$ 没有预设关系。

取局部标架 $\{e_a\}$ 与余标架 $\{\theta^a\}$。定义联络系数
\begin{equation*}\nabla_{e_b}e_c=\Gamma^a{}_{bc}e_a,
\end{equation*}
并把同一信息组织成联络一形式
\begin{equation*}\nabla e_c=\omega^a{}_c\otimes e_a,
\qquad
\omega^a{}_c=\Gamma^a{}_{bc}\theta^b.
\end{equation*}
于是对 $Y=Y^ce_c$，
\[
\nabla_XY
=X(Y^a)e_a+Y^c\omega^a{}_c(X)e_a.
\]
这显示了联络一形式的角色：它补偿标架自身随位置的变化，使“分量求导 + 基底求导”重新组合成几何向量。

联络自然延伸到任意张量。对一形式 $\alpha$，定义
\[
(\nabla_X\alpha)(Y)
:=X[\alpha(Y)]-\alpha(\nabla_XY),
\]
再要求与张量积和缩并相容，就唯一确定一般张量的协变导数。对含标架指标的微分形式还需要把普通外微分升级为协变外微分。若 $\eta=\eta^ae_a$ 是向量值 $p$-形式，定义
\begin{equation}\label{eq:math-dg-exterior-covariant-derivative}
(D\eta)^a
:=\dd\eta^a+\omega^a{}_b\wedge\eta^b.
\end{equation}
若指标为协变指标，相应联络项带负号；多个指标时对每个指标分别作用。$D$ 是普通 $\dd$ 在“带联络的向量丛值形式”上的延伸，后面的 Bianchi 恒等式都用这一定义。

联络系数不是张量。若换标架 $e'=e\Lambda$，其中 $\Lambda:U\to GL(n,\RR)$，则
\begin{equation}\label{eq:math-dg-connection-gauge-transform}
\boxed{\omega'=\Lambda^{-1}\omega\Lambda+\Lambda^{-1}\dd\Lambda.}
\end{equation}
非齐次项 $\Lambda^{-1}\dd\Lambda$ 正是“联络不是张量”的精确含义。

\begin{derivation}[联络一形式的换标架律]
联络由 $\nabla e_c=e_a\otimes\omega^a{}_c$ 定义。换标架 $e'=e\Lambda$ 后，对 $\nabla e'_b$ 用 Leibniz 展开再换回新基，读出新系数即得非齐次变换律。

Leibniz 展开。设新标架 $e'_b=e_c\Lambda^c{}_b$，其中 $\Lambda:U\to GL(n,\RR)$ 光滑。由联络对被求导向量满足 Leibniz 规则，对 $\nabla e'_b=\nabla(e_c\Lambda^c{}_b)$ 逐项展开，
\begin{align*}
\nabla e'_b
&=\nabla(e_c\Lambda^c{}_b)\\
&=(\nabla e_c)\Lambda^c{}_b+e_c\otimes\dd\Lambda^c{}_b.
\end{align*}
第一项是"基向量被联络作用"再乘系数，第二项是"系数被外微分"乘基向量——这正是联络作为方向导数的 Leibniz 性质。

代入旧联络。由旧标架下联络一形式的定义 $\nabla e_c=e_a\otimes\omega^a{}_c$，代回第一项，
\begin{equation*}
(\nabla e_c)\Lambda^c{}_b
=e_a\otimes\omega^a{}_c\Lambda^c{}_b.
\end{equation*}

合并两项。把指标 $c$ 在第一项中作哑指标，$a$ 作自由指标，合并两项
\begin{align*}
\nabla e'_b
&=e_a\otimes\omega^a{}_c\Lambda^c{}_b
 +e_c\otimes\dd\Lambda^c{}_b\\
&=e_a\otimes\bigl(\omega^a{}_c\Lambda^c{}_b+\dd\Lambda^a{}_b\bigr),
\end{align*}
其中第二行把第二项的哑指标 $c$ 改写为 $a$。

换回新基。为读出新标架下的联络一形式 $\omega'$，需把基向量 $e_a$ 换回 $e'_r$。由 $e'_b=e_c\Lambda^c{}_b$ 取逆得 $e_a=e'_r(\Lambda^{-1})^r{}_a$，代回
\begin{align*}
\nabla e'_b
&=e'_r(\Lambda^{-1})^r{}_a\otimes
 \bigl(\omega^a{}_c\Lambda^c{}_b+\dd\Lambda^a{}_b\bigr)\\
&=e'_r\otimes
 \Bigl[
  (\Lambda^{-1})^r{}_a\omega^a{}_c\Lambda^c{}_b
  +(\Lambda^{-1})^r{}_a\dd\Lambda^a{}_b
 \Bigr].
\end{align*}

读出系数。与新标架下定义 $\nabla e'_b=e'_r\otimes\omega'^r{}_b$ 比较基向量 $e'_r$ 的系数，得
\begin{equation*}
\omega'^r{}_b
=(\Lambda^{-1})^r{}_a\omega^a{}_c\Lambda^c{}_b
+(\Lambda^{-1})^r{}_a\dd\Lambda^a{}_b,
\end{equation*}
即\eqnrefs{eq:math-dg-connection-gauge-transform}。第一项是齐次的共轭转动，第二项是非齐次的纯规范补偿。

典型例子。
\begin{itemize}
\item 正交换标架 $\Lambda\in SO(n)$：$\Lambda^{-1}=\Lambda^T$，$\omega'=\Lambda^T\omega\Lambda+\Lambda^T\dd\Lambda$。第二项取值于 $\mathfrak{so}(n)$（反对称），保持度量相容性。
\item 局部 Lorentz 标架（tetrad）变换 $\Lambda\in SO(1,3)$：$\omega'=\Lambda^{-1}\omega\Lambda+\Lambda^{-1}\dd\Lambda$，正是广义相对论中 spin connection 在局部 Lorentz 变换下的规律。
\item 复标架变换 $\Lambda\in U(N)$：$\omega'=\Lambda^\dagger\omega\Lambda+\Lambda^\dagger\dd\Lambda$，对应 $U(N)$ 规范势的规范变换。
\end{itemize}

非齐次项 $\Lambda^{-1}\dd\Lambda$ 是"联络不是张量"的精确含义。在广义相对论中，它保证协变导数 $\nabla_\mu V^\nu=\partial_\mu V^\nu+\Gamma^\nu{}_{\mu\rho}V^\rho$ 在坐标变换下张量性成立——$\Gamma$ 的非齐次变换正好抵消 $\partial_\mu V^\nu$ 中多出的二阶项。在规范理论中，它使规范势 $A_\mu$ 在规范变换下获得 $\partial_\mu\chi$ 补偿项，从而协变导数 $D_\mu\psi$ 齐次变换。
\end{derivation}

向量丛值形式上的曲率最简洁地由 $D^2$ 定义。对 $E$-值 $p$-形式 $\eta=\eta^ae_a$，令
\[
\Omega:=\dd\omega+\omega\wedge\omega.
\]
则有
\begin{equation}\label{eq:math-dg-covariant-derivative-square}
\boxed{D^2\eta=\Omega\wedge\eta.}
\end{equation}
这条公式先于切丛的指标表达成立；对 $E=TM$ 才把 $\Omega$ 解释为 Riemann 曲率二形式。

\begin{derivation}[协变外微分平方等于曲率作用]
由协变外微分的定义 $D\eta=\dd\eta+\omega\wedge\eta$，其中 $\omega$ 是取值于 $\End E$ 的联络一形式、$\eta$ 是 $E$-值 $p$-形式。再次作用 $D$，先把 $D\eta$ 视作新的 $E$-值 $(p+1)$-形式代入同一公式，
\begin{equation*}
D^2\eta
=D(D\eta)
=\dd(D\eta)+\omega\wedge(D\eta).
\end{equation*}
把 $D\eta=\dd\eta+\omega\wedge\eta$ 代入两处，
\begin{equation*}
D^2\eta
=\dd(\dd\eta+\omega\wedge\eta)
 +\omega\wedge(\dd\eta+\omega\wedge\eta).
\end{equation*}
由外微分的线性性拆开，
\begin{equation*}
D^2\eta
=\dd^2\eta+\dd(\omega\wedge\eta)
 +\omega\wedge\dd\eta+\omega\wedge(\omega\wedge\eta).
\end{equation*}
由 $\dd^2=0$ 第一项消失。对第二项用 graded Leibniz 规则 $\dd(\alpha\wedge\beta)=\dd\alpha\wedge\beta+(-1)^{|\alpha|}\alpha\wedge\dd\beta$，其中 $|\omega|=1$，
\begin{equation*}
\dd(\omega\wedge\eta)
=\dd\omega\wedge\eta+(-1)^1\omega\wedge\dd\eta
=\dd\omega\wedge\eta-\omega\wedge\dd\eta.
\end{equation*}
代回，
\begin{equation*}
D^2\eta
=\dd\omega\wedge\eta-\omega\wedge\dd\eta
 +\omega\wedge\dd\eta+\omega\wedge\omega\wedge\eta.
\end{equation*}
中间两项 $-\omega\wedge\dd\eta+\omega\wedge\dd\eta$ 严格相消，只剩
\begin{equation*}
D^2\eta
=\dd\omega\wedge\eta+\omega\wedge\omega\wedge\eta
=(\dd\omega+\omega\wedge\omega)\wedge\eta.
\end{equation*}
由曲率二形式定义 $\Omega:=\dd\omega+\omega\wedge\omega$，代回得
\begin{equation*}
D^2\eta=\Omega\wedge\eta,
\end{equation*}
即\eqnrefs{eq:math-dg-covariant-derivative-square}。这也解释了为何曲率必然以 $\dd\omega+\omega\wedge\omega$ 而不是单独 $\dd\omega$ 出现。
\end{derivation}

联络本身不是向量空间中的元素：两个联络可以相减，但单独一个联络不能在不选参考点的情况下被称为“零”。设 $\nabla$ 与 $\widetilde\nabla$ 是同一向量丛 $E$ 上的两个联络，定义
\[
\mathcal A_Xs
:=(\widetilde\nabla_X-\nabla_X)s.
\]
Leibniz 项精确相消，因此 $\mathcal A_Xs$ 对 $X$ 与 $s$ 都是 $C^\infty(M)$-线性的。故
\begin{equation*}\boxed{
\widetilde\nabla-\nabla
\in\Omega^1(M;\End E).
}\end{equation*}
反过来，任意 $A\in\Omega^1(M;\End E)$ 都定义新联络
\[
(\nabla+A)_Xs:=\nabla_Xs+A(X)s.
\]
所以所有联络组成一个以 $\Omega^1(M;\End E)$ 为模型向量空间的 affine space。

在局部标架中，若参考联络形式为 $\omega$，则新联络形式只是
\[
\widetilde\omega=\omega+A.
\]
但曲率不是线性变化。由
\[
F_\nabla=\dd\omega+\omega\wedge\omega
\]
直接展开得
\begin{equation}
\boxed{
F_{\nabla+A}
=F_\nabla+D_\nabla A+A\wedge A,
}
\label{eq:math-dg-curvature-finite-perturbation}
\end{equation}
其中
\[
D_\nabla A
:=\dd A+\omega\wedge A+A\wedge\omega
\]
是 $\End(E)$-值一形式上的协变外微分。若取一参数族 $\nabla_t=\nabla+tA$，则一阶变分为
\begin{equation*}\left.\frac{\dd}{\dd t}\right|_{t=0}F_{\nabla_t}
=D_\nabla A.\end{equation*}
Chern--Weil transgression、Yang--Mills 线性化以及后面度量变分中的 $\partial_tR$ 都是这一一般公式的不同表现。

对切丛 affine connection，差张量可写成
\[
A(X,Y):=(\widetilde\nabla_X-\nabla_X)Y
\in\Gamma(TM).
\]
挠率的变化特别简单：
\begin{equation*}\boxed{
\widetilde T(X,Y)
=T(X,Y)+A(X,Y)-A(Y,X).
}\end{equation*}
所以给定一个联络后，只需调整差张量的反对称部分就能改变 torsion；其对称部分不影响 torsion。曲率则满足有限差分公式
\begin{align}
\widetilde R(X,Y)Z
={}&R(X,Y)Z
+(\nabla_XA)(Y,Z)-(\nabla_YA)(X,Z)\notag\\
&+A(X,A(Y,Z))-A(Y,A(X,Z))\notag\\
&+A(T(X,Y),Z),
\label{eq:math-dg-curvature-connection-change}
\end{align}
其中最后一项反映参考联络可能本身有 torsion；若 $T=0$ 就消失。

\begin{derivation}[曲率有限扰动为什么必然含线性项与二次项，即\eqnrefs{eq:math-dg-curvature-finite-perturbation}]
曲率是联络的二次函数（含 $\nabla^2$ 和 $\nabla\wedge\nabla$），故联络的有限扰动 $A$ 产生线性项（$D_\nabla A$）和二次项（$A\wedge A$），结构完全由交换子代数决定。

分解新联络。把新联络写成
\begin{equation*}
 \widetilde\nabla_X=\nabla_X+A_X,
 \qquad
 A_X:=A(X,\cdot).
\end{equation*}
$A$ 是 $\Omega^1(M;\operatorname{End}E)$ 中任一元素，$\widetilde\nabla$ 与 $\nabla$ 同为 $E$ 上联络当且仅当它们的差是张量一形式。

代入曲率定义。曲率是交换子减联络沿括号方向的作用：
\begin{equation*}
 \widetilde R(X,Y)
 =[\widetilde\nabla_X,\widetilde\nabla_Y]-\widetilde\nabla_{[X,Y]}
 =[\nabla_X+A_X,\nabla_Y+A_Y]
 -(\nabla_{[X,Y]}+A_{[X,Y]}).
\end{equation*}

展开交换子。按双线性展开：
\begin{itemize}
\item 纯 $\nabla$ 部分：$[\nabla_X,\nabla_Y]-\nabla_{[X,Y]}=R(X,Y)$，即旧曲率。
\item 一次 $A$ 的交叉项：$\nabla_XA_Y-\nabla_YA_X-A_{[X,Y]}$，组合成 covariant exterior derivative $(D_\nabla A)(X,Y)$，即 $A$ 沿旧联络的协变外微分。
\item 二次项：$[A_X,A_Y]=A_XA_Y-A_YA_X$，即 $A\wedge A$。
\end{itemize}
逐项核对：$[\nabla_X,A_Y]=\nabla_XA_Y-A_Y\nabla_X$ 中 $-A_Y\nabla_X$ 与 $[A_X,\nabla_Y]=A_X\nabla_Y-\nabla_YA_X$ 中的 $A_X\nabla_Y$ 不直接相消，但与 $\nabla_XA_Y-\nabla_YA_X$ 合并后正好是协变外微分 $D_\nabla A$ 的定义。

抽象形式。因此
\begin{equation*}
 \widetilde R=R+D_\nabla A+A\wedge A,
\end{equation*}
与\eqnrefs{eq:math-dg-curvature-finite-perturbation}完全相同。曲率对联络"线性化 + 自相互作用"的结构因此不是坐标巧合，而是 connection affine geometry 的必然结果：线性项 $D_\nabla A$ 是一阶变分（用于 Chern--Weil 理论和指标定理），二次项 $A\wedge A$ 是非 Abel 自相互作用（与 Yang--Mills 方程的非线性性同源）。

线性化与 Chern--Weil。取 $A_t=tA$，$\widetilde R_t=R+tD_\nabla A+t^2 A\wedge A$。对 $\operatorname{tr}(\widetilde R_t^k)$ 在 $t=0$ 求导，只有线性项 $D_\nabla A$ 贡献（二次项 $t^2$ 在 $t=0$ 导数为零）。这正是一阶变分控制 Chern--Weil transgression 的原因：特征类的变分只看 $D_\nabla A$，非 Abel 自相互作用 $A\wedge A$ 在一阶变分中不可见，但在二阶变分与 Yang--Mills 作用量中起决定作用。

Yang--Mills 作用量的二次展开。Yang--Mills 作用量 $S=\int\operatorname{tr}(F\wedge *F)$ 在 $\widetilde\nabla=\nabla+A$ 下展开为
\begin{equation*}
 S(\widetilde\nabla)=S(\nabla)+2\int\operatorname{tr}(D_\nabla A\wedge *F)+\int\operatorname{tr}((D_\nabla A+A\wedge A)\wedge *(D_\nabla A+A\wedge A)).
\end{equation*}
第一阶变分给 Yang--Mills 方程 $D_\nabla *F=0$；二阶变分含 $A\wedge A$ 的二次与三次项，是规范场自能的非线性来源。Abel 情形 $A\wedge A=0$，作用量退化为 Maxwell 二次型，变分方程线性。

典型例子。
\begin{itemize}
\item $U(1)$ 规范理论（Abel）：$A\wedge A=0$（交换子为零），$\widetilde F=F+\dd A$。电磁场强度在规范势平移下只加 $\dd A$，没有自相互作用——这正是 Maxwell 方程线性的根源。
\item $SU(2)$ Yang--Mills：$\widetilde F=F+D_\nabla A+A\wedge A$，$A\wedge A\neq0$ 产生非 Abel 自相互作用。Yang--Mills 方程 $D_\mu F^{\mu\nu}=0$ 展开后含 $A_\mu A_\nu A^\nu$ 三次项，是非线性性的来源。
\item Riemann 几何：取 $\nabla$ 为 Levi--Civita 联络、$A$ 为 $(1,2)$ 张量，$\widetilde R=R+\nabla A+A\circ A$ 给出两个度量相容联络的曲率差。若 $\widetilde\nabla$ 也无挠，$A$ 对下指标对称，二次项控制如 Rauch 比较定理中的曲率比较。
\end{itemize}

这一展开是规范理论全部相互作用结构的微缩版：线性项 $D_\nabla A$ 是"外源激发"，二次项 $A\wedge A$ 是"自相互作用"。在 Yang--Mills 中，规范场既被外源激发又自我激发，正是非 Abel 规范理论区别于 Maxwell 理论的本质。在广义相对论中，Einstein 张量 $G_{\mu\nu}$ 含 Ricci 曲率，而 Ricci 在联络扰动下的二次项对应引力场的自相互作用——这是 Einstein 方程非线性的几何根源。
\end{derivation}

若 $E$ 带 fiber metric $h$，保持 $h$ 的联络差张量必须取值于 $h$-反自伴 endomorphisms；在正交标架下就是 $\mathfrak{so}(r)$-值一形式。特别地，对 Riemannian tangent bundle，metric-compatible connections 构成由 $\Omega^1(M;\mathfrak{so}(n))$ 平移得到的 affine 子空间。下一章再加入 torsion-free 条件后，这个 affine freedom 被完全固定，得到唯一 Levi--Civita 联络。

\section{挠率、曲率与 Cartan 结构方程}
联络给定后，先比较 $\nabla_XY-\nabla_YX$ 与光滑结构已经给出的 $[X,Y]$。它们的差定义挠率
\begin{equation*}T(X,Y)
:=\nabla_XY-\nabla_YX-[X,Y].
\end{equation*}
逐项检查可知对 $X,Y$ 都是 $C^\infty(M)$-线性的，因此 $T$ 是真正的 $(1,2)$ 张量。写
\[
T(X,Y)=T^a(X,Y)e_a
\]
定义挠率二形式 $T^a$。

曲率则比较沿 $X,Y$ 两个方向连续协变微分的交换子，并扣除方向场本身的不交换：
\begin{equation*}R(X,Y)Z
:=\nabla_X\nabla_YZ-\nabla_Y\nabla_XZ-\nabla_{[X,Y]}Z.
\end{equation*}
它对 $X,Y,Z$ 都是函数线性的，因此是 $(1,3)$ 张量。定义曲率二形式 $\Omega^a{}_b$ 为
\[
R(X,Y)e_b=\Omega^a{}_b(X,Y)e_a.
\]
若
\[
R(e_c,e_d)e_b=R^a{}_{bcd}e_a,
\]
则
\[
\Omega^a{}_b=\frac12R^a{}_{bcd}\theta^c\wedge\theta^d.
\]

\begin{derivation}[Cartan 第一、第二结构方程]
先从挠率定义出发。对任意 $X,Y$，
\[
\begin{aligned}
T^a(X,Y)
={}&\theta^a(\nabla_XY-\nabla_YX-[X,Y])\\
={}&X[\theta^a(Y)]-Y[\theta^a(X)]-\theta^a([X,Y])\\
&+\omega^a{}_b(X)\theta^b(Y)
-\omega^a{}_b(Y)\theta^b(X).
\end{aligned}
\]
第一行中前三项恰是 $\dd\theta^a(X,Y)$，后二项恰是
$(\omega^a{}_b\wedge\theta^b)(X,Y)$。故
\begin{equation}\label{eq:math-dg-synthesis-cartan-first}
\boxed{T^a=\dd\theta^a+\omega^a{}_b\wedge\theta^b.}
\end{equation}

再对基向量 $e_b$ 展开曲率定义：
\[
\begin{aligned}
R(X,Y)e_b
={}&\nabla_X\bigl(\omega^c{}_b(Y)e_c\bigr)
-\nabla_Y\bigl(\omega^c{}_b(X)e_c\bigr)
-\omega^c{}_b([X,Y])e_c\\
={}&\Bigl(
X[\omega^a{}_b(Y)]-Y[\omega^a{}_b(X)]-\omega^a{}_b([X,Y])
\Bigr)e_a\\
&+\Bigl(
\omega^c{}_b(Y)\omega^a{}_c(X)
-\omega^c{}_b(X)\omega^a{}_c(Y)
\Bigr)e_a.
\end{aligned}
\]
把第一括号认作 $\dd\omega^a{}_b(X,Y)$，第二括号认作
$(\omega^a{}_c\wedge\omega^c{}_b)(X,Y)$，得到
\begin{equation}\label{eq:math-dg-synthesis-cartan-second}
\boxed{
\Omega^a{}_b
=\dd\omega^a{}_b+\omega^a{}_c\wedge\omega^c{}_b.}
\end{equation}

几何意义。第一结构方程 $T^a=\dd\theta^a+\omega^a{}_b\wedge\theta^b$ 说挠率是「外微分加联络修正」；第二结构方程 $\Omega^a{}_b=\dd\omega^a{}_b+\omega^a{}_c\wedge\omega^c{}_b$ 说曲率是联络的「外微分加自乘」。二者合起来把挠率和曲率完全用标架 $\{\theta^a\}$ 和联络 $\{\omega^a{}_b\}$ 表达，无需坐标。对 Riemannian 流形（挠率 $T=0$），第一方程显式解出 $\omega^a{}_b$，将其表为 $\dd\theta^a$ 的函数，第二方程则给出曲率——这是 Cartan 活动标架法计算曲率的标准流程。

二维球面。取球面 $\mathbb S^2$ 的活动标架 $\theta^1=R\,\dd\theta$、$\theta^2=R\sin\theta\,\dd\varphi$。由 $T=0$ 解第一方程得 $\omega^1{}_2=-\cos\theta\,\dd\varphi$。代入第二方程：$\Omega^1{}_2=\dd\omega^1{}_2=\sin\theta\,\dd\theta\wedge\dd\varphi=(1/R^2)\,\theta^1\wedge\theta^2$，立即读出截面曲率 $K=1/R^2$。整个计算无需 Christoffel 符号。
\end{derivation}

Cartan 结构方程由此写成外微分形式：挠率 $T^a=\dd\theta^a+\omega^a{}_b\wedge\theta^b$（式~\eqnrefs{eq:math-dg-synthesis-cartan-first}），曲率 $\Omega^a{}_b=\dd\omega^a{}_b+\omega^a{}_c\wedge\omega^c{}_b$（式~\eqnrefs{eq:math-dg-synthesis-cartan-second}）。

在一般标架中，把 $\omega^a{}_c=\Gamma^a{}_{bc}\theta^b$ 与
$[e_b,e_c]=C^a{}_{bc}e_a$ 代入第一结构方程，挠率分量为
\begin{equation*}T^a{}_{bc}
=\Gamma^a{}_{bc}-\Gamma^a{}_{cb}-C^a{}_{bc}.
\end{equation*}
因此“联络系数下标对称”只在坐标基 $C^a{}_{bc}=0$ 中才与无挠等价。非坐标基中若忘掉结构函数，会直接得到错误结论。

在 principal bundle 语言中，曲率形式为
\begin{equation*}\mathbf F
=\dd\boldsymbol\omega
+\frac12[\boldsymbol\omega\wedge\boldsymbol\omega],\end{equation*}
拉回局部截面后成为
\begin{equation*}F=\dd A+A\wedge A.\end{equation*}
换 gauge 时没有非齐次项：
\begin{equation*}F'=g^{-1}Fg.\end{equation*}
因此曲率是 tensorial object；它真正测量联络本身，而不是局部 frame 的选择。

若 $F=0$，称联络 flat。单连通坐标域上，平坦性正是局部存在 gauge $g$ 使
\begin{equation*}A=g^{-1}\dd g\end{equation*}
的可积条件；再作 gauge transformation $g^{-1}$ 就可令局部 connection form 消失。于是
\[
F=0
\quad\Longleftrightarrow\quad
A\text{ 局部可 gauge 到 }0
\]
在单连通局部成立。但全局平坦联络仍可能具有非平凡 holonomy，因为 noncontractible loops 能保存全局拓扑信息。

\begin{derivation}[曲率的齐次 gauge 变换，即 $F'=g^{-1}Fg$]
出发点。gauge 变换下 connection form 的变换律为
\begin{equation*}
 A'=g^{-1}Ag+g^{-1}\dd g.
\end{equation*}
记 $\eta=g^{-1}\dd g$（纯 gauge），则 $A'=g^{-1}Ag+\eta$。目标是对 $F'=\dd A'+A'\wedge A'$ 直接计算并证明 $F'=g^{-1}Fg$。

Maurer--Cartan 恒等式。由 $\eta=g^{-1}\dd g$ 是纯 gauge，Maurer--Cartan 方程给出
\begin{equation*}
 \dd\eta+\eta\wedge\eta=0.
\end{equation*}
这是 $\eta$ 的"平坦性"条件，几何意义是 $g$ 是全局定义的光滑函数，其微分可积。

再计算 $g^{-1}Ag$ 的外微分。由 $\dd(g^{-1})=-\eta\,g^{-1}$（对 $g\cdot g^{-1}=I$ 求微分），
\begin{equation*}
 \dd(g^{-1}Ag)
 =g^{-1}(\dd A)g
 -\eta\wedge g^{-1}Ag
 -g^{-1}Ag\wedge\eta.
\end{equation*}
后两项是 $\eta$ 与 $g^{-1}Ag$ 的交叉楔积，来自 $\dd(g^{-1})$ 与 $\dd g$ 上的 Leibniz 规则。

代入 $F'=\dd A'+A'\wedge A'$。展开后收集所有项：
\begin{itemize}
\item $g^{-1}(\dd A)g$：来自 $\dd(g^{-1}Ag)$。
\item $g^{-1}(A\wedge A)g$：来自 $(g^{-1}Ag)\wedge(g^{-1}Ag)$。
\item $\dd\eta$：来自 $\dd\eta$。
\item $\eta\wedge\eta$：来自 $\eta\wedge\eta$。
\item 交叉项 $\eta\wedge g^{-1}Ag$ 和 $g^{-1}Ag\wedge\eta$：来自 $\dd(g^{-1}Ag)$ 展开的后两项和 $A'\wedge A'$ 的展开。
\end{itemize}
具体地，$A'\wedge A'=(g^{-1}Ag+\eta)\wedge(g^{-1}Ag+\eta)$ 展开为四项：$(g^{-1}Ag)\wedge(g^{-1}Ag)$、$(g^{-1}Ag)\wedge\eta$、$\eta\wedge(g^{-1}Ag)$、$\eta\wedge\eta$。其中两个交叉项与 $\dd(g^{-1}Ag)$ 中的两个交叉项符号相反，精确配对抵消。

用 Maurer--Cartan 收尾。由 Maurer--Cartan 恒等式 $\dd\eta=-\eta\wedge\eta$，以及交叉项精确配对抵消，只剩
\begin{equation*}
 F'=g^{-1}(\dd A+A\wedge A)g
 =g^{-1}Fg.
\end{equation*}
交叉项的相消可显式核对：$\dd(g^{-1}Ag)$ 给 $-\eta\wedge g^{-1}Ag-g^{-1}Ag\wedge\eta$，$A'\wedge A'$ 给 $+g^{-1}Ag\wedge\eta+\eta\wedge g^{-1}Ag$，二者符号相反；$\dd\eta$ 与 $\eta\wedge\eta$ 则由 Maurer--Cartan 恒等式消去。

典型例子。
\begin{itemize}
\item $U(1)$ 规范：$g=e^{\ii\chi}$，$F'=F$（Abel 群共轭平凡）。电磁场强度 $F_{\mu\nu}=\partial_\mu A_\nu-\partial_\nu A_\mu$ 在规范变换 $A_\mu\mapsto A_\mu+\partial_\mu\chi$ 下不变，正是 Maxwell 理论规范不变性的几何根源。
\item $SU(N)$ Yang--Mills：$F'=g^{-1}Fg$，规范场强做共轭变换。$\operatorname{tr}(F\wedge F)$ 因此规范不变，是 Chern--Weil 理论构造第二 Chern 类 $c_2(E)$ 的基础。
\item $SO(n)$ 标架丛：$\Omega'=g^{-1}\Omega g$ 给 Riemann 曲率二形式的齐次变换律，与联络一形式的非齐次律对比鲜明。
\item 仿射联络：$\Gamma$ 在坐标变换下非齐次（含二阶导数项），但 Riemann 曲率张量 $R^\rho{}_{\sigma\mu\nu}$ 齐次变换，是真正的张量。
\item 瞬子解：$SU(2)$ BPST 瞬子 $A_\mu=(1-f)\,g^{-1}\partial_\mu g$ 中 $F=g^{-1}F_0g$ 共轭不变，$\operatorname{tr}(F\wedge F)=\operatorname{tr}(F_0\wedge F_0)$ 给瞬子数 $k=1$。
\end{itemize}

这就是"connection 非张量而 curvature 张量"的矩阵形式证明：$A$ 在 gauge 变换下有非齐次项 $g^{-1}\dd g$（仿射变换），但 $F$ 只做共轭变换（齐次变换），因此曲率是真正的张量。在 Yang--Mills 理论中，作用量 $\int\operatorname{tr}(F\wedge *F)$ 规范不变正是源于 $F$ 的齐次变换律与 trace 的循环性。在 Chern--Weil 理论中，$\operatorname{tr}(F^k)$ 规范不变且闭，定义了不依赖联络选择的拓扑不变量。这一"非齐次 $\to$ 齐次"的升级是规范理论全部构造性结论（守恒流、规范不变作用量、拓扑荷）的共同根源。
\end{derivation}

曲率二形式的换标架律则是齐次的：
\begin{equation*}\Omega'=\Lambda^{-1}\Omega\Lambda.
\end{equation*}
这表明曲率本身是几何张量，虽然表示它的联络一形式不是。

\section{Bianchi 恒等式}
Cartan 方程并不是任意二形式方程。因为 $T$ 与 $\Omega$ 都来自同一个联络，它们自动满足两条微分一致性条件。

\begin{derivation}[第一与第二 Bianchi 恒等式]
由第一结构方程，
\[
T^a=\dd\theta^a+\omega^a{}_b\wedge\theta^b.
\]
用\eqnrefs{eq:math-dg-exterior-covariant-derivative}计算
\[
\begin{aligned}
D T^a
={}&\dd T^a+\omega^a{}_c\wedge T^c\\
={}&\dd\omega^a{}_b\wedge\theta^b
-\omega^a{}_b\wedge\dd\theta^b
+\omega^a{}_c\wedge\dd\theta^c
+\omega^a{}_c\wedge\omega^c{}_b\wedge\theta^b.
\end{aligned}
\]
中间两项在重命名哑指标后相消，于是
\[
D T^a
=\bigl(\dd\omega^a{}_b+\omega^a{}_c\wedge\omega^c{}_b\bigr)\wedge\theta^b
=\Omega^a{}_b\wedge\theta^b.
\]
因此
\begin{equation}\label{eq:math-dg-bianchi-first}
\boxed{D T^a=\Omega^a{}_b\wedge\theta^b.}
\end{equation}

对第二结构方程作协变外微分。对一个 $\End(TM)$-值二形式，
\[
(D\Omega)^a{}_b
=\dd\Omega^a{}_b
+\omega^a{}_c\wedge\Omega^c{}_b
-\Omega^a{}_c\wedge\omega^c{}_b.
\]
代入 $\Omega=\dd\omega+\omega\wedge\omega$，并使用 $\dd^2=0$，
\[
\begin{aligned}
D\Omega
={}&\dd(\omega\wedge\omega)
+\omega\wedge(\dd\omega+\omega\wedge\omega)
-(\dd\omega+\omega\wedge\omega)\wedge\omega\\
={}&(\dd\omega\wedge\omega-\omega\wedge\dd\omega)
+\omega\wedge\dd\omega+\omega\wedge\omega\wedge\omega\\
&-\dd\omega\wedge\omega-\omega\wedge\omega\wedge\omega\\
={}&0.
\end{aligned}
\]
故
\begin{equation}\label{eq:math-dg-bianchi-second}
\boxed{D\Omega^a{}_b=0.}
\end{equation}
两式也可合写为
\begin{equation}\label{eq:math-dg-synthesis-bianchi}
D T^a=\Omega^a{}_b\wedge\theta^b,
\qquad
D\Omega^a{}_b=0.
\end{equation}

意义。Bianchi 恒等式是曲率的可积条件：$D\Omega=0$ 说曲率 2-形式在外协变微分下闭，类比于 $\dd^2=0$。这是 Chern--Weil 理论的出发点——$\operatorname{tr}(\Omega^k)$ 在 $D$ 下闭，因此定义了不依赖联络选取的示性类。实例：二维曲面上 $D\Omega=0$ 退化为 $\dd K\,\operatorname{vol}=0$ 在常曲率情形自动成立；四维 Yang--Mills 理论中 $D\Omega=0$ 是第二 Bianchi 恒等式的规范理论形式，与 Bianchi 恒等式 $D\Omega=0$ 共同保证 $\nabla_\mu G^{\mu\nu}=0$（Einstein 方程的协变守恒）。
\end{derivation}

由此得到结构方程的 Bianchi 恒等式：$DT^a=\Omega^a{}_b\wedge\theta^b$ 与 $D\Omega^a{}_b=0$（式~\eqnrefs{eq:math-dg-bianchi-first}、\eqnrefs{eq:math-dg-bianchi-second}、\eqnrefs{eq:math-dg-synthesis-bianchi}）。

第二 Bianchi 恒等式的一个深层后果是：对任意 Ad-invariant polynomial $P$ on $\mathfrak g$，把曲率代入所得微分形式
\begin{equation*}P(F)\in\Omega^{2k}(M)\end{equation*}
自动闭合。最常见的矩阵表示例子是
\[
\operatorname{tr}(F^k).
\]
由于 trace 消去 commutator，
\[
\dd\operatorname{tr}(F^k)
=\operatorname{tr}(D F^k)
=k\operatorname{tr}((DF)F^{k-1})=0.
\]
因此它定义 de Rham cohomology class。

更重要的是，这个 cohomology class 不依赖具体 connection。令 $A_t$ 是一族联络，$F_t$ 对应曲率。则
\begin{equation}
\frac{\dd}{\dd t}\operatorname{tr}(F_t^k)
=k\,\dd\operatorname{tr}
\left(\dot A_t\wedge F_t^{k-1}\right).
\label{eq:chern-weil-transgression}
\end{equation}
右端是 exact form，所以 $[\operatorname{tr}(F_t^k)]$ 与 $t$ 无关。曲率虽然依赖联络，某些由曲率构造的上同调类却只依赖 bundle topology；这就是 Chern--Weil 理论的核心机制。

\begin{derivation}[Chern--Weil class 为何与 connection 无关]
对一族联络 $A_t$ 与曲率 $F_t$，把 $\operatorname{tr}(F_t^k)$ 对 $t$ 求导。利用 $\dot F_t=D_t\dot A_t$、Bianchi 恒等式 $D_tF_t=0$ 与 trace 的循环性，把导数写成 exact form，从而 de Rham 上同调类与 $t$ 无关。

曲率对联络的变分。对
\begin{equation*}
F_t=\dd A_t+A_t\wedge A_t
\end{equation*}
求导（$\dot A_t:=\frac{\dd A_t}{\dd t}$）：
\begin{equation*}
\dot F_t
=\dd\dot A_t+A_t\wedge\dot A_t+
\dot A_t\wedge A_t
=D_t\dot A_t.
\end{equation*}
最后等式是协变外微分的定义 $D_t\dot A_t=\dd\dot A_t+[A_t,\dot A_t]$，其中 $[A_t,\dot A_t]=A_t\wedge\dot A_t+\dot A_t\wedge A_t$（一形式楔积带符号合并）。

对 $\operatorname{tr}(F_t^k)$ 求导。由 Leibniz 规则，
\begin{equation*}
\frac{\dd}{\dd t}\operatorname{tr}(F_t^k)
=k\,\operatorname{tr}(\dot F_t\wedge F_t^{k-1})
=k\,\operatorname{tr}(D_t\dot A_t\wedge F_t^{k-1}).
\end{equation*}
因子 $k$ 来自 $F_t^k=F_t\wedge\cdots\wedge F_t$ 中 $\dot F_t$ 可落在 $k$ 个位置之一，trace 的循环性使这 $k$ 项相等。

利用协变 Leibniz 规则和 Bianchi 恒等式，可把 $D_t$ 从 $\dot A_t$ 移到 $F_t^{k-1}$：
\begin{equation*}
D_t(\dot A_t\wedge F_t^{k-1})
=D_t\dot A_t\wedge F_t^{k-1}
+\dot A_t\wedge D_t(F_t^{k-1}),
\end{equation*}
以及第二 Bianchi 恒等式 $D_tF_t=0$（从而 $D_t(F_t^{k-1})=0$），得
\begin{equation*}
D_t(\dot A_t\wedge F_t^{k-1})
=D_t\dot A_t\wedge F_t^{k-1}.
\end{equation*}
这一步是关键：Bianchi 恒等式让 $D_t$ 在乘积上"穿过"所有 $F_t$ 因子，仿佛它们是"协变常量"。

trace 把 $D_t$ 退化为 $\dd$。对 trace 而言，$D_t$ 与普通 $\dd$ 的差是 commutator $[A_t,\cdot]$，其 trace 为零（trace 的循环性 $\operatorname{tr}([X,Y])=\operatorname{tr}(XY)-\operatorname{tr}(YX)=0$）。因此
\begin{equation*}
\operatorname{tr}(D_t(\dot A_t\wedge F_t^{k-1}))
=\dd\operatorname{tr}(\dot A_t\wedge F_t^{k-1}),
\end{equation*}
得到\eqnrefs{eq:chern-weil-transgression}。

积分得 de Rham 类不变。积分 $t\in[0,1]$ 后，
\begin{equation*}
\operatorname{tr}(F_1^k)-\operatorname{tr}(F_0^k)
=k\,\dd\int_0^1\operatorname{tr}(\dot A_t\wedge F_t^{k-1})\,\dd t,
\end{equation*}
两组 connection 产生的 characteristic forms 相差 exact form，因此代表同一 de Rham class $[\operatorname{tr}(F^k)]\in H_{\mathrm{dR}}^{2k}(M)$。

典型例子。
\begin{itemize}
\item $U(1)$ 线丛：$k=1$，$\operatorname{tr}(F)=F$ 是闭二形式，$[F/2\pi]\in H^2_{\mathrm{dR}}(M)$ 是第一 Chern 类 $c_1(E)$。两个不同联络的 $F$ 相差 exact form，但积分 $\frac{1}{2\pi}\int_M F$ 给出 Chern 数（如磁通量子化）。
\item $SU(N)$ 丛：$k=2$，$\operatorname{tr}(F\wedge F)$ 是闭四形式，$[\operatorname{tr}(F\wedge F)/8\pi^2]$ 是第二 Chern 类 $c_2(E)$。Yang--Mills 瞬子数 $k=\int\operatorname{tr}(F\wedge F)/8\pi^2$ 是拓扑不变量，不依赖规范选择。
\end{itemize}

更一般结论。Chern--Weil 理论给出从 invariant polynomials on $\mathfrak g$ 到 $H_{\mathrm{dR}}^*(M)$ 的同态：每个 $G$-不变多项式 $P$ 给出闭形式 $P(F)$，其 de Rham 类只依赖丛 $E$ 的拓扑。结合 Atiyah--Singer 指标定理，这些特征类积分给出椭圆算子的解析指标——如 Dirac 算子的 $\hat A$ 类、符号丛的 Chern 字符。这是"曲率局部数据决定全局拓扑"的精确机制。
\end{derivation}

对复向量丛，适当归一化的
\[
\det\left(I+\frac{\ii}{2\pi}F\right)
=1+c_1(E)+c_2(E)+\cdots
\]
给出 Chern classes 的 de Rham representatives；对实定向丛则由 curvature 的 invariant polynomials 构造 Pontryagin classes 与 Euler class。这里不展开完整特征类理论，但应保留逻辑顺序：
\[
\text{bundle}
\to\text{connection}
\to\text{curvature}
\to\text{closed invariant forms}
\to\text{topological cohomology classes}.
\]
这也是为什么曲率既是局部平行移动缺陷，又能够携带全局拓扑信息。

第一 Bianchi 恒等式把挠率与曲率联系起来；第二 Bianchi 恒等式表达曲率自身的协变闭合。它们都不是运动方程，而是由联络定义强制产生的恒等式。

\section{平行移动与曲率的几何意义}
沿光滑曲线 $\gamma(t)$，向量场 $V(t)\in T_{\gamma(t)}M$ 若满足
\begin{equation*}\nabla_{\dot\gamma}V=0,
\end{equation*}
则称 $V$ 沿 $\gamma$ 平行。写 $V=V^ae_a$ 后，这是线性常微分方程
\[
\frac{\dd V^a}{\dd t}
+\omega^a{}_b(\dot\gamma)V^b=0.
\]
因此给定联络以后，任意曲线都诱导一个线性平行移动映射。


局部标架中，平行移动方程可写成矩阵 ODE
\begin{equation}
\dot V(t)+A_tV(t)=0,
\qquad
A_t:=\omega(\dot\gamma(t)).
\label{eq:parallel-matrix-ode}
\end{equation}
若不同时刻的 $A_t$ 彼此对易，则
\[
V(t)=\exp\left(-\int_0^tA_\tau\,\dd\tau\right)V(0).
\]
一般情况下必须保留时间次序，写成 path-ordered exponential
\begin{equation*}U_\gamma(t,0)
=\mathcal P\exp\left(-\int_0^tA_\tau\,\dd\tau\right).\end{equation*}
其 Dyson series 为
\[
I+
\sum_{k=1}^\infty(-1)^k
\int_{0<t_k<\cdots<t_1<t}
A_{t_1}\cdots A_{t_k}
\,\dd t_k\cdots\dd t_1.
\]
所以 path ordering 不是记号装饰，而是非交换 connection matrices 的精确解法。

若 $\gamma$ 是闭路，$U_\gamma$ 给出 holonomy transformation。对面积为 $\epsilon^2$ 的微小平行四边形，其 holonomy 有展开
\begin{equation*}U_{\partial\Sigma}
=I-
F(X,Y)\epsilon^2+O(\epsilon^3),\end{equation*}
符号取决于平行移动方程和闭路定向约定。核心内容是二阶系数恰为 curvature。后续 holonomy 章节会把这个局部关系提升为 holonomy algebra 与曲率生成元的全局定理。

\begin{derivation}[path-ordered exponential 的迭代构造]
把\eqnrefs{eq:parallel-matrix-ode} 即 $\dot V(t)=-A_tV(t)$、$V(0)$ 给定，从 $0$ 到 $t$ 积分得积分方程
\begin{equation*}
V(t)-V(0)=-\int_0^tA_{t_1}V(t_1)\,\dd t_1,
\end{equation*}
即
\begin{equation*}
V(t)=V(0)-\int_0^tA_{t_1}V(t_1)\,\dd t_1.
\end{equation*}
把右端 $V(t_1)$ 用同一积分方程代入，$V(t_1)=V(0)-\int_0^{t_1}A_{t_2}V(t_2)\,\dd t_2$，得
\begin{align*}
V(t)
&=V(0)-\int_0^tA_{t_1}
 \Bigl[V(0)-\int_0^{t_1}A_{t_2}V(t_2)\,\dd t_2\Bigr]\dd t_1\\
&=V(0)-\int_0^tA_{t_1}\dd t_1\,V(0)
 +\int_0^t\int_0^{t_1}A_{t_1}A_{t_2}V(t_2)\,\dd t_2\dd t_1\\
&=\Bigl[I-\int_0^tA_{t_1}\dd t_1\Bigr]V(0)
 +\int_0^t\int_0^{t_1}A_{t_1}A_{t_2}V(t_2)\,\dd t_2\dd t_1.
\end{align*}
对余项中的 $V(t_2)$ 再代入一次，$V(t_2)=V(0)-\int_0^{t_2}A_{t_3}V(t_3)\,\dd t_3$，
\begin{align*}
\int_0^t\int_0^{t_1}A_{t_1}A_{t_2}V(t_2)\,\dd t_2\dd t_1
&=\int_0^t\int_0^{t_1}A_{t_1}A_{t_2}\dd t_2\dd t_1\,V(0)\\
&\quad-\int_0^t\int_0^{t_1}\int_0^{t_2}
 A_{t_1}A_{t_2}A_{t_3}V(t_3)\,\dd t_3\dd t_2\dd t_1.
\end{align*}
合并得第三次迭代
\begin{align*}
V(t)
&=\Bigl[I-\int_0^tA_{t_1}\dd t_1
 +\int_0^t\int_0^{t_1}A_{t_1}A_{t_2}\dd t_2\dd t_1\Bigr]V(0)\\
&\quad-\int_0^t\int_0^{t_1}\int_0^{t_2}
 A_{t_1}A_{t_2}A_{t_3}V(t_3)\,\dd t_3\dd t_2\dd t_1.
\end{align*}
由归纳法，第 $k$ 次迭代把余项写成在有序区域
\begin{equation*}
0<t_k<\cdots<t_2<t_1<t
\end{equation*}
上对 $A_{t_1}\cdots A_{t_k}V(t_k)$ 的 $k$ 重积分，前 $k$ 项之和为
\begin{equation*}
\sum_{j=0}^{k}(-1)^j
\int_{0<t_j<\cdots<t_1<t}
A_{t_1}\cdots A_{t_j}\,\dd t_j\cdots\dd t_1\,V(0),
\end{equation*}
其中 $j=0$ 项理解为 $I\,V(0)$。在 $A_t$ 关于 $t$ 一致有界的前提下余项趋于零，从而产生 Dyson series
\begin{equation*}
V(t)=\Bigl[I+\sum_{k=1}^\infty(-1)^k
\int_{0<t_k<\cdots<t_1<t}
A_{t_1}\cdots A_{t_k}\,\dd t_k\cdots\dd t_1\Bigr]V(0).
\end{equation*}
若所有 $A_t$ 对易，积分区域的 $k!$ 个排列给相同 integrand，于是有序单纯形积分等于立方体积分的 $1/k!$，级数便退化为普通指数。这解释了什么时候可以安全地把 $\mathcal P$ 去掉。
\end{derivation}

曲率的局部意义可以从无穷小闭合回路看出。若先沿 $X$ 再沿 $Y$ 平行移动，并与相反次序比较，扣除底点回路本身由 $[X,Y]$ 造成的偏移以后，最低非平凡差异正由 $R(X,Y)$ 给出。换言之，$R=0$ 意味着局部平行移动对足够小的可缩闭路没有二阶 holonomy；$R\neq0$ 则说明不同路径的比较无法彼此兼容。
